\section{Introduction}

Phenotypic drug discovery asks how a chemical or genetic perturbation changes the state of a cell, at a scale that target-specific assays cannot reach. Bringing one drug to approval costs on the order of \$2\,B and takes about a decade~\cite{wouters2020rdcost}, while roughly nine of ten compounds that enter clinical trials fail~\cite{sun2022clinicaltrials}.

High-content microscopy is one way to get that scale. The Cell Painting assay~\cite{bray2016cellpainting, seal2024cellpainting} stains cells with five fluorescent dyes that mark DNA, endoplasmic reticulum, RNA and nucleoli, actin and Golgi, and mitochondria, and images them as a five-channel picture of cellular state. A perturbation changes that picture, and the change carries information about the perturbation's mechanism. The CPJUMP1 benchmark~\cite{chandrasekaran2024jump} turns this into a retrieval problem over about 3 million fields of view: given images of cells treated with 303 compounds, 160 CRISPR knockouts and 176 ORF overexpressions, can a model match perturbations that act on the same target?

Existing methods each cover part of this problem, and none covers all of it. The CellProfiler pipeline~\cite{carpenter2006cellprofiler} measures handcrafted features and recovers 5 to 25\% of the expected sister-perturbation matches on CPJUMP1~\cite{chandrasekaran2024jump}. CLOOME~\cite{sanchez2023cloome} and MolPhenix~\cite{fradkin2024molphenix} learn contrastive image encoders against molecular fingerprints, and so handle compounds only. CellCLIP~\cite{lu2025cellclip} replaces fingerprints with text prompts, which admits genetic perturbations, but runs on a 1.48B-parameter DINOv2-g stack and its cross-class retrieval trails CellProfiler. CWA-MSN~\cite{huang2025cwamsn} attacks batch effects directly and beats CellCLIP on gene--gene retrieval with 22M parameters, but has no text branch and cannot be queried in language.

MorphoCLIP combines text supervision, joint training over compounds and genes, and a batch-correction option in one model. It follows CLIP~\cite{radford2021clip}: an image encoder and a text encoder are trained so that a Cell Painting well and the description of its perturbation land close together in a shared 512-dimensional space. Both backbones are frozen and their outputs are cached. The trainable components are a one-layer transformer over channel tokens, two projection heads and one logit-scale scalar, about 14M parameters in total, compared with CellCLIP's 1.48B total parameters.

The paper makes four contributions.
\begin{enumerate}
    \item A text-supervised model over cached DINOv3 channel tokens~\cite{simeoni2025dinov3}, trained jointly on compounds, CRISPR knockouts and ORF overexpressions with two frozen backbones.
    \item Three additions measured separately against a shared control: an image--image replicate loss, gene-aware soft labels, and a per-plate offset relative to replicate plates of the same condition. The replicate loss increases deterministic replicability mAP on every benchmark track; retrieval differences and the other additions remain unresolved at one seed.
    \item Perturbation-level retrieval with analytic random baselines and evaluation on the standard CPJUMP1 benchmark. Separate stochastic benchmark runs show numerically comparable compound replicability fractions for MorphoCLIP and the released CellCLIP checkpoint, without supporting a ranking. The four-track mean compound target-matching fraction lies within the published CellProfiler range, although one track is above it. Gene--compound fraction retrieved remains at or near zero.
    \item A documented evaluation pitfall: duplicating each text once per replicate well creates tied blocks of candidates, caps Recall@5 and Recall@10 near Recall@1, and invalidates a shuffled-well baseline for image-to-text retrieval.
\end{enumerate}
